\section{Two Authorizations That Do Not Know About Each Other}
\label{sec:ch15-two-authorizations}
\index{HotChocolate!two authorization mechanisms|(}%

HotChocolate 16.6.0 gives you two ways to say that a field is protected. They
come from different packages, they are enforced in different processes, and
nothing in the product connects them. I knew both existed. What I did not know,
until I put one on each of two fields in the same service and looked at what
came out, is that neither one leaves any trace of the other.

\subsection*{The one the service enforces}

\index{Authorize@\code{[Authorize]}!enforced in process}%
\code{[Authorize]} comes from \code{HotChocolate.Authorization}, and once
ASP.NET Core is in the picture the builder call comes from
\code{HotChocolate.AspNetCore.Authorization}. Accounts puts it on the one root
field it has:

\begin{minted}[fontsize=\footnotesize,breaklines]{csharp}
[Authorize]
public static Task<Customer?> GetCustomerByIdAsync(
    [ID] Guid id,
    AccountsService accounts,
    CancellationToken cancellationToken)
    => accounts.GetCustomerByIdAsync(id, cancellationToken);
\end{minted}

The attribute takes a policy name, a role list and an \code{ApplyPolicy} that
decides whether the rule runs before the resolver, after it, or during
validation. Accounts uses none of those and gets the bare meaning: there has to
be somebody logged in.

That rule is evaluated by this process, against the principal this process
built out of the token this process validated. The router is not involved and
does not know it exists.

\subsection*{The one the router enforces}

\index{RequiresScopes@\code{[RequiresScopes]}!published to the router}%
\code{[RequiresScopes]} comes from the federation package, along with
\code{[Authenticated]} and \code{[Policy]}. It goes on the field it protects,
which for a positional record means the property:

\begin{minted}[fontsize=\footnotesize,breaklines]{csharp}
[property: RequiresScopes(["read:pii"])]
string Email)
\end{minted}

Its entire effect is on the schema the service publishes. Accounts does not
check scopes, cannot check scopes, and answers this field to an anonymous
caller that asks it directly.

One detail here is worth carrying, because it is in no documentation I could
find and it reads backwards. The attribute takes a flat \code{string[]}; the
directive it produces takes \code{[[Scope!]!]!}. The attribute is declared
\code{AllowMultiple = true}. So one attribute is one AND-group, every scope in
it required together, and a second attribute stacked beside it is the OR. Write
\code{[RequiresScopes(["a", "b"])]} expecting \enquote{either} and you have
written \enquote{both}.

\subsection*{What each one publishes}

\index{schema export@\code{schema export}!what the attributes emit}%
Both attributes are on fields of one service, so one \code{schema export} shows
what each contributes. Against the committed schema from
chapter~\ref{ch:real-time-in-a-federated-world}, four things changed:

\begin{graphqlsdl}
import: ["@key", "@requiresScopes", "@tag", "FieldSet"]

customerById(id: ID!): Customer @authorize

email: String! @requiresScopes(scopes: [["read:pii"]])

enum ApplyPolicy {
  BEFORE_RESOLVER
  AFTER_RESOLVER
  VALIDATION
}
\end{graphqlsdl}

The doc comments the exporter writes above each of those enum members are cut
here for room; nothing else in the four is abridged.

The first is HotChocolate adding \code{@requiresScopes} to the \code{@link}
import list on its own, which is tidy of it. The third is the directive doing
exactly what it is for, and it is the only one of the four that was intended.

The second is the surprise, and I had believed the opposite on good grounds.
The type behind the attribute calls \code{.Internal()} on the directive, and
an internal directive is one that stays out of introspection. I read that and
concluded that \code{[Authorize]} publishes nothing. It publishes
\code{@authorize} on the field and a full declaration of the directive
underneath, and what a composer reads is the SDL a subgraph serves from
\code{\_service}, not its introspection result. The lesson
is one this book has learned before in a different costume: what a service
publishes is a thing to look at, not a thing to reason about.

The fourth is the one I would not have predicted at all. \code{enum
ApplyPolicy} appeared in all seven schemas, including the six services that use
no authorization attribute anywhere. It is the argument type of
\code{@authorize}, and it arrives with the \code{.AddAuthorization()} call that
\code{AddMosaicSubgraph} makes for every subgraph. Six services gained a public
type because a shared method gained a line.

\subsection*{And what survives composition}

\index{ApplyPolicy@\code{ApplyPolicy}!reaches the client schema}%
Composing those seven schemas sorts the four into two piles.
\code{@requiresScopes} survives, which
section~\ref{sec:ch15-what-the-composer-checks} takes apart properly.
\code{@authorize} is dropped: it appears nowhere in the composed client schema,
and the composer says nothing about having dropped it.

\code{ApplyPolicy} is not dropped. It is a type of Mosaic's graph now,
introspectable by every client, describing when a policy runs inside a
HotChocolate service. Chapter~\ref{ch:composition} found the same shape in
\code{scalar FieldSet}, which reaches the client schema because the composer
keeps imported named types, and which chapter~\ref{ch:security-hardening}
inherited as a question about what a production graph does with it. This is a
second instance and it is worse, because \code{FieldSet} at least belongs to
federation. \code{ApplyPolicy} is an implementation detail of one server
library, in the public schema of a graph, and no part of the toolchain thinks
that is worth a word.

Figure~\ref{fig:ch15-two-authorizations} is the whole arrangement on one page:
two attributes on two fields of one service, two paths that never meet, and a
composer in the middle that keeps one of them and drops the other without
comment.

\begin{figure}[htbp]
  \centering
  % Two authorization mechanisms on one service, and nothing between them.
% Referenced from chapter 15. Bare tikzpicture; the figure environment,
% caption and label live at the call site.
%
% The point of the picture is the vertical divider. It runs the whole height
% of the figure and breaks only once, exactly at the composer, which is the
% only place the two attributes ever appear together. Nothing crosses from
% one side to the other anywhere else: the left lane is checked and consumed
% entirely inside Accounts, the right lane is never read by Accounts at all,
% and the composer is what decides which of the two the router ever hears
% about.
\begin{tikzpicture}[
    font=\small,
    tp/.style={draw, rounded corners=2pt, minimum width=42mm,
               minimum height=15mm, align=center, font=\scriptsize},
    dirbox/.style={draw, rounded corners=2pt, minimum width=40mm,
                   minimum height=13mm, align=center, font=\scriptsize},
    composer/.style={draw, rounded corners=2pt, minimum width=46mm,
                     minimum height=13mm, align=center, font=\scriptsize,
                     fill=black!8},
    dropped/.style={draw, dashed, rounded corners=2pt, minimum width=40mm,
                    minimum height=20mm, align=center, font=\scriptsize,
                    gray},
    kept/.style={draw, rounded corners=2pt, minimum width=46mm,
                minimum height=24mm, align=center, font=\scriptsize,
                fill=black!8},
    flow/.style={-{Stealth[length=2mm]}},
    divider/.style={gray, dashed},
    note/.style={font=\scriptsize, gray, align=center, inner sep=1pt},
    closing/.style={font=\scriptsize\itshape, align=center}
  ]

  % -- one subgraph, two fields -----------------------------------------------
  \node[note, gray, anchor=south, font=\small] at (6.5,1.65) {\code{accounts}};
  \draw[rounded corners=2pt] (0.3,-2.4) rectangle (12.7,1.5);

  \node[tp] (f1) at (3.2,0.5)
    {\code{Query.customerById}\\\code{[Authorize]}\\\code{HotChocolate.Authorization}};
  \node[tp] (f2) at (9.8,0.5)
    {\code{Customer.email}\\\code{[RequiresScopes]}\\\code{HotChocolate.ApolloFederation}};

  \node[note, anchor=north] at (3.2,-0.85)
    {evaluated here,\\against the principal\\built from the token};
  \node[note, anchor=north] at (9.8,-0.85) {Accounts never reads it};

  % -- both compile to SDL, published either way -------------------------------
  \node[dirbox] (d1) at (3.2,-3.5) {\code{@authorize}\\on \code{customerById}};
  \node[dirbox] (d2) at (9.8,-3.5)
    {\code{@requiresScopes}\\\code{(scopes: [["read:pii"]])}};

  \draw[flow] (f1.south) -- (d1.north);
  \draw[flow] (f2.south) -- (d2.north);

  \node[note, anchor=east] at (2.6,-2.6) {compiles to;\\published in\\the SDL};
  \node[note, anchor=west] at (10.4,-2.6) {compiles to;\\published in\\the SDL};

  % -- the only place the two lanes meet ---------------------------------------
  \node[composer] (comp) at (6.5,-5.6)
    {\code{wgc router compose}\\the composer};

  \draw[flow] (d1) -- (comp);
  \draw[flow] (d2) -- (comp);

  \draw[divider] (6.5,1.75) -- (6.5,-5.0);
  \draw[divider] (6.5,-6.2) -- (6.5,-9.0);

  % -- and the two fates ------------------------------------------------------
  \node[dropped] (drop) at (3.2,-8.0)
    {\textbf{silently dropped}\\by the composer\\reaches the router\\as nothing};
  \node[kept] (keep) at (9.8,-8.0)
    {\textbf{survives composition}\\\code{fieldConfigurations}\\entry in the
     router's\\execution config;\\the router refuses};

  \draw[flow] (6.5,-6.2) -- (6.5,-6.6) -- (3.2,-6.6) -- (drop.north);
  \draw[flow] (6.5,-6.2) -- (6.5,-6.6) -- (9.8,-6.6) -- (keep.north);

  \draw[gray, thick] (1.8,-9.1) -- (4.6,-9.1);

  \node[closing, text width=124mm] at (6.5,-10.0)
    {Nothing on the left produces the directive on the right, and nothing on
     the right knows the left exists. The two attributes share one service and
     meet, once, inside the composer, which keeps one and throws the other
     away.};

\end{tikzpicture}

  \caption{Two authorization mechanisms in one subgraph, and the gap between
  them. \code{[Authorize]} is enforced by Accounts and publishes an
  \code{@authorize} directive the composer discards, so nothing about it
  reaches the router. \code{[RequiresScopes]} is never read by Accounts and
  survives composition into the router's execution config, where it is
  enforced. Nothing in the product generates either one from the other, and
  nothing warns when only one of them is present.}
  \label{fig:ch15-two-authorizations}
\end{figure}

So the state of things after one attribute each: the service enforces a rule
the router cannot see, the router will enforce a rule the service does not
check, and the only thing both of them agree to publish is an enum neither of
them meant.

\medskip

None of that matters yet, because nothing in this graph has ever accepted a
token. Getting one accepted is the next section, and it is harder than it
looks.

\index{HotChocolate!two authorization mechanisms|)}%
